\documentclass[11pt, a4paper]{article}
\usepackage[a4paper, margin=2cm]{geometry}
\usepackage{fontspec}
\usepackage[english, bidi=basic, provide=*]{babel}
\babelprovide[import, onchar=ids fonts]{english}
\babelfont{rm}{Noto Sans}
\usepackage{amsmath}
\usepackage{amssymb}
\usepackage{booktabs}
\usepackage{enumitem}
\usepackage[table]{xcolor}
\usepackage{hyperref}

\definecolor{primary}{RGB}{15, 76, 129}
\definecolor{secondary}{RGB}{44, 62, 80}

\hypersetup{
    colorlinks=true,
    linkcolor=primary,
    urlcolor=primary
}

\begin{document}

\begin{center}
    {\Large \textbf{\color{primary} NEET BIOLOGY: HIGH-DIFFICULTY NCERT PRACTICE BANK}} \\[6pt]
    {\large \textbf{\color{secondary} Subtopics: Mitochondria, Plastids \& Ribosomes}} \\[4pt]
    {\small \textit{45 Strictly NCERT-Based Conceptual, Statement, Assertion-Reason \& Match-the-Column Questions}}
\end{center}

\vspace{4pt}
\hrule
\vspace{10pt}

\section*{\color{primary} Part I: Multiple Choice Questions (45 Questions)}

\begin{enumerate}[label=\textbf{\arabic*.}, leftmargin=*, itemsep=8pt]

    \item \textbf{Statement I:} Mitochondria are typically sausage-shaped or cylindrical organelles having a diameter of $0.2\text{--}1.0\,\mu\text{m}$ (average $0.5\,\mu\text{m}$) and length of $1.0\text{--}4.1\,\mu\text{m}$.\\
    \textbf{Statement II:} Mitochondria are easily visible under the light microscope without requiring any specific vital staining.\\
    In light of the above statements, choose the correct option:
    \begin{enumerate}[label=(\alph*)]
        \item Both Statement I and Statement II are correct
        \item Both Statement I and Statement II are incorrect
        \item Statement I is correct but Statement II is incorrect
        \item Statement I is incorrect but Statement II is correct
    \end{enumerate}

    \item Consider the following statements regarding the number, shape, and size of mitochondria in eukaryotic cells:
    \begin{enumerate}[label=(\roman*)]
        \item The number of mitochondria per cell is constant across all cell types within a given organism.
        \item Considerable degree of variability is observed in the shape and size of mitochondria.
        \item Cells with higher physiological activity possess a greater number of mitochondria per cell.
        \item Each mitochondrion is a double membrane-bound structure with the inner compartment filled with a dense homogeneous substance called the matrix.
    \end{enumerate}
    How many of the above statements are \textbf{correct}?
    \begin{enumerate}[label=(\alph*)]
        \item One
        \item Two
        \item Three
        \item Four
    \end{enumerate}

    \item \textbf{Assertion (A):} The inner mitochondrial membrane forms a number of infoldings called the cristae towards the matrix.\\
    \textbf{Reason (R):} The cristae significantly decrease the surface area of the inner membrane so that enzymes of the electron transport chain remain tightly packed in a small space.\\
    Choose the correct option:
    \begin{enumerate}[label=(\alph*)]
        \item Both (A) and (R) are true and (R) is the correct explanation of (A)
        \item Both (A) and (R) are true but (R) is not the correct explanation of (A)
        \item (A) is true but (R) is false
        \item Both (A) and (R) are false
    \end{enumerate}

    \item Which of the following features is \textbf{incorrectly} attributed to the mitochondrial matrix as per NCERT?
    \begin{enumerate}[label=(\alph*)]
        \item It possesses a single circular DNA molecule
        \item It contains a few RNA molecules
        \item It houses $80\text{S}$ ribosomes for synthesizing mitochondrial proteins
        \item It contains components required for the synthesis of proteins
    \end{enumerate}

    \item Match Column I with Column II regarding the structural organization of mitochondria:
    \begin{center}
    \begin{tabular}{@{}ll@{}}
        \toprule
        \textbf{Column I (Component)} & \textbf{Column II (Characteristic Feature)} \\
        \midrule
        A. Outer membrane & 1. Divides the lumen distinctly into two aqueous compartments \\
        B. Inner membrane & 2. Forms the continuous limiting boundary of the organelle \\
        C. Cristae & 3. Site of cellular respiration producing cellular energy in the form of ATP \\
        D. Mitochondria & 4. Infoldings towards the matrix that increase surface area \\
        \bottomrule
    \end{tabular}
    \end{center}
    Choose the correct match:
    \begin{enumerate}[label=(\alph*)]
        \item A-2, B-1, C-4, D-3
        \item A-1, B-2, C-4, D-3
        \item A-2, B-4, C-1, D-3
        \item A-3, B-1, C-4, D-2
    \end{enumerate}

    \item How do mitochondria multiply inside a eukaryotic cell?
    \begin{enumerate}[label=(\alph*)]
        \item By budding from the outer nuclear membrane
        \item By binary fission
        \item By multiple fission during S-phase only
        \item By de novo assembly from endoplasmic reticulum vesicles
    \end{enumerate}

    \item \textbf{Statement I:} The outer and inner membranes of a mitochondrion have their own specific enzymes associated with the mitochondrial function.\\
    \textbf{Statement II:} Mitochondria are known as the "power houses" of the cell because they are the sites of anaerobic fermentation.\\
    Choose the correct option:
    \begin{enumerate}[label=(\alph*)]
        \item Both Statement I and Statement II are correct
        \item Both Statement I and Statement II are incorrect
        \item Statement I is correct but Statement II is incorrect
        \item Statement I is incorrect but Statement II is correct
    \end{enumerate}

    \item Read the following statements about plastids:
    \begin{enumerate}[label=(\roman*)]
        \item Plastids are found in all plant cells and in euglenoids.
        \item They are easily observed under the microscope as they are large.
        \item They bear some specific pigments, thus imparting specific colours to the plants.
        \item Based on the type of pigments, plastids can be classified into chloroplasts, chromoplasts, and leucoplasts.
    \end{enumerate}
    Which of the above statements are \textbf{correct}?
    \begin{enumerate}[label=(\alph*)]
        \item (i) and (ii) only
        \item (ii), (iii), and (iv) only
        \item (i), (iii), and (iv) only
        \item (i), (ii), (iii), and (iv)
    \end{enumerate}

    \item Which of the following pigments are present in chromoplasts to give parts of the plant a yellow, orange, or red colour?
    \begin{enumerate}[label=(\alph*)]
        \item Water-soluble anthocyanins and chlorophyll $b$
        \item Fat-soluble carotenoid pigments like carotene, xanthophylls, and others
        \item Water-soluble phycobilins and chlorophyll $a$
        \item Fat-soluble anthoxanthins and haemoglobin
    \end{enumerate}

    \item Match Column I with Column II regarding the types of leucoplasts and their stored nutrients:
    \begin{center}
    \begin{tabular}{@{}ll@{}}
        \toprule
        \textbf{Column I (Type of Leucoplast)} & \textbf{Column II (Stored Nutrient / Example)} \\
        \midrule
        A. Amyloplasts & 1. Store proteins \\
        B. Elaioplasts & 2. Store carbohydrates (starch), e.g., potato \\
        C. Aleuroplasts & 3. Store oils and fats \\
        D. Leucoplasts & 4. Colourless plastids of varied shapes and sizes \\
        \bottomrule
    \end{tabular}
    \end{center}
    Choose the correct option:
    \begin{enumerate}[label=(\alph*)]
        \item A-2, B-3, C-1, D-4
        \item A-3, B-2, C-1, D-4
        \item A-2, B-1, C-3, D-4
        \item A-1, B-3, C-2, D-4
    \end{enumerate}

    \item \textbf{Assertion (A):} In chloroplasts, the inner membrane is relatively less permeable than the outer membrane.\\
    \textbf{Reason (R):} Like mitochondria, the chloroplasts are also double membrane-bound organelles.\\
    Choose the correct option:
    \begin{enumerate}[label=(\alph*)]
        \item Both (A) and (R) are true and (R) is the correct explanation of (A)
        \item Both (A) and (R) are true but (R) is not the correct explanation of (A)
        \item (A) is true but (R) is false
        \item Both (A) and (R) are false
    \end{enumerate}

    \item Identify the \textbf{incorrect} statement regarding the dimensions and number of chloroplasts in plant cells:
    \begin{enumerate}[label=(\alph*)]
        \item Majority of the chloroplasts of the green plants are found in the mesophyll cells of the leaves
        \item They are lens-shaped, oval, spherical, discoid, or even ribbon-like organelles
        \item Their length ranges from $5\text{--}10\,\mu\text{m}$ and width ranges from $2\text{--}4\,\mu\text{m}$
        \item \textit{Chlamydomonas} contains $20\text{--}40$ chloroplasts per cell, whereas a mesophyll cell contains only $1$ chloroplast per cell
    \end{enumerate}

    \item Consider the following statements about the internal organization of a chloroplast:
    \begin{enumerate}[label=(\roman*)]
        \item The space limited by the inner membrane of the chloroplast is called the stroma.
        \item A number of organised flattened membranous sacs called the thylakoids are present in the stroma.
        \item Thylakoids are arranged in stacks like the piles of coins called grana (singular: granum) or the intergranal thylakoids.
        \item Flat membranous tubules called the stroma lamellae connect the thylakoids of the different grana.
    \end{enumerate}
    How many of the above statements are \textbf{correct}?
    \begin{enumerate}[label=(\alph*)]
        \item One
        \item Two
        \item Three
        \item Four
    \end{enumerate}

    \item \textbf{Statement I:} The membrane of the thylakoids encloses a space called a lumen.\\
    \textbf{Statement II:} Chlorophyll pigments are located exclusively in the stroma of the chloroplast.\\
    Choose the correct option:
    \begin{enumerate}[label=(\alph*)]
        \item Both Statement I and Statement II are correct
        \item Both Statement I and Statement II are incorrect
        \item Statement I is correct but Statement II is incorrect
        \item Statement I is incorrect but Statement II is correct
    \end{enumerate}

    \item Which of the following components are found within the stroma of a chloroplast?
    \begin{enumerate}[label=(\alph*)]
        \item Enzymes required for the synthesis of carbohydrates and proteins only
        \item Small, double-stranded circular DNA molecules and $70\text{S}$ ribosomes only
        \item Enzymes required for the synthesis of carbohydrates and proteins, small double-stranded circular DNA molecules, and $70\text{S}$ ribosomes
        \item Single-stranded linear DNA molecules, $80\text{S}$ ribosomes, and hydrolytic acid hydrolases
    \end{enumerate}

    \item Why are mitochondria and chloroplasts excluded from the endomembrane system of a eukaryotic cell?
    \begin{enumerate}[label=(\alph*)]
        \item Because they are single membrane-bound organelles
        \item Because their functions are not coordinated with the endoplasmic reticulum, Golgi complex, lysosomes, and vacuoles
        \item Because they lack ribosomes and genetic material
        \item Because they are present only in prokaryotic cells
    \end{enumerate}

    \item Compare the ribosomes of the eukaryotic cytoplasm with those found inside chloroplasts and mitochondria:
    \begin{enumerate}[label=(\alph*)]
        \item Cytoplasmic ribosomes are $70\text{S}$, whereas organelle ribosomes are $80\text{S}$
        \item Cytoplasmic ribosomes are $80\text{S}$, whereas organelle ribosomes are $70\text{S}$
        \item Both cytoplasmic and organelle ribosomes are $80\text{S}$
        \item Both cytoplasmic and organelle ribosomes are $70\text{S}$
    \end{enumerate}

    \item \textbf{Assertion (A):} Ribosomes are non-membrane bound organelles found in all cells—both eukaryotic as well as prokaryotic.\\
    \textbf{Reason (R):} Within a eukaryotic cell, ribosomes are found not only in the cytoplasm but also within the two organelles—chloroplasts (in plants) and mitochondria—and on rough ER.\\
    Choose the correct option:
    \begin{enumerate}[label=(\alph*)]
        \item Both (A) and (R) are true and (R) is the correct explanation of (A)
        \item Both (A) and (R) are true but (R) is not the correct explanation of (A)
        \item (A) is true but (R) is false
        \item Both (A) and (R) are false
    \end{enumerate}

    \item Who first observed ribosomes under the electron microscope as dense particles in the year 1953?
    \begin{enumerate}[label=(\alph*)]
        \item Camillo Golgi
        \item George Palade
        \item Robert Brown
        \item Singer and Nicolson
    \end{enumerate}

    \item \textbf{Statement I:} Ribosomes are composed of ribonucleic acid (RNA) and proteins and are not surrounded by any membrane.\\
    \textbf{Statement II:} In the sedimentation coefficient notation ($70\text{S}$ and $80\text{S}$), 'S' stands for Svedberg's unit, which directly measures the molecular weight of the ribosome.\\
    Choose the correct option:
    \begin{enumerate}[label=(\alph*)]
        \item Both Statement I and Statement II are correct
        \item Both Statement I and Statement II are incorrect
        \item Statement I is correct but Statement II is incorrect
        \item Statement I is incorrect but Statement II is correct
    \end{enumerate}

    \item Identify the correct combination of subunits for eukaryotic and prokaryotic ribosomes respectively:
    \begin{enumerate}[label=(\alph*)]
        \item $60\text{S} + 40\text{S}$ (for $80\text{S}$) and $50\text{S} + 30\text{S}$ (for $70\text{S}$)
        \item $50\text{S} + 30\text{S}$ (for $80\text{S}$) and $40\text{S} + 30\text{S}$ (for $70\text{S}$)
        \item $60\text{S} + 20\text{S}$ (for $80\text{S}$) and $50\text{S} + 20\text{S}$ (for $70\text{S}$)
        \item $40\text{S} + 40\text{S}$ (for $80\text{S}$) and $35\text{S} + 35\text{S}$ (for $70\text{S}$)
    \end{enumerate}

    \item In prokaryotes, ribosomes are associated with the plasma membrane of the cell and measure about:
    \begin{enumerate}[label=(\alph*)]
        \item $5\text{ nm}$ by $8\text{ nm}$
        \item $15\text{ nm}$ by $20\text{ nm}$
        \item $50\text{ nm}$ by $80\text{ nm}$
        \item $0.2\,\mu\text{m}$ by $1.0\,\mu\text{m}$
    \end{enumerate}

    \item \textbf{Statement I:} Several ribosomes may attach to a single mRNA and form a chain called polyribosomes or polysome.\\
    \textbf{Statement II:} The ribosomes of a polysome translate the mRNA into proteins.\\
    Choose the correct option:
    \begin{enumerate}[label=(\alph*)]
        \item Both Statement I and Statement II are correct
        \item Both Statement I and Statement II are incorrect
        \item Statement I is correct but Statement II is incorrect
        \item Statement I is incorrect but Statement II is correct
    \end{enumerate}

    \item Which of the following features is \textbf{shared} by both mitochondria and chloroplasts?
    \begin{enumerate}[label=(\alph*)]
        \item Presence of thylakoids and grana
        \item Presence of cristae on the inner membrane
        \item Presence of double membrane, circular DNA, and $70\text{S}$ ribosomes
        \item Presence of fat-soluble carotenoid pigments
    \end{enumerate}

    \item Read the following statements about plastids and identify the \textbf{incorrect} one:
    \begin{enumerate}[label=(\alph*)]
        \item Leucoplasts are the colourless plastids of varied shapes and sizes with stored nutrients
        \item In chromoplasts, chlorophyll is completely absent, and only water-soluble pigments are present
        \item Chloroplasts contain chlorophyll and carotenoid pigments which are responsible for trapping light energy essential for photosynthesis
        \item Plastids are absent in prokaryotic cells and animal cells
    \end{enumerate}

    \item Match Column I with Column II based on the exact dimensions and characteristics given in NCERT:
    \begin{center}
    \begin{tabular}{@{}ll@{}}
        \toprule
        \textbf{Column I (Structure)} & \textbf{Column II (NCERT Dimension / Feature)} \\
        \midrule
        A. Mitochondrion diameter & 1. $5\text{--}10\,\mu\text{m}$ \\
        B. Chloroplast length & 2. $0.2\text{--}1.0\,\mu\text{m}$ (average $0.5\,\mu\text{m}$) \\
        C. Prokaryotic ribosome & 3. $2\text{--}4\,\mu\text{m}$ \\
        D. Chloroplast width & 4. $15\text{ nm}$ by $20\text{ nm}$ \\
        \bottomrule
    \end{tabular}
    \end{center}
    Choose the correct option:
    \begin{enumerate}[label=(\alph*)]
        \item A-2, B-1, C-4, D-3
        \item A-1, B-2, C-4, D-3
        \item A-2, B-3, C-4, D-1
        \item A-4, B-1, C-2, D-3
    \end{enumerate}

    \item \textbf{Assertion (A):} In a chloroplast, the light-independent (dark) reactions of carbohydrate synthesis take place in the stroma.\\
    \textbf{Reason (R):} The stroma of the chloroplast contains the enzymes required for the synthesis of carbohydrates and proteins.\\
    Choose the correct option:
    \begin{enumerate}[label=(\alph*)]
        \item Both (A) and (R) are true and (R) is the correct explanation of (A)
        \item Both (A) and (R) are true but (R) is not the correct explanation of (A)
        \item (A) is true but (R) is false
        \item Both (A) and (R) are false
    \end{enumerate}

    \item Consider the following statements regarding the aqueous compartments of a mitochondrion:
    \begin{enumerate}[label=(\roman*)]
        \item The inner membrane divides the mitochondrial lumen distinctly into two aqueous compartments: the outer compartment and the inner compartment.
        \item The outer compartment is filled with a dense homogeneous substance called the matrix.
        \item The outer membrane forms the continuous limiting boundary of the organelle.
        \item Both membranes have identical enzymatic compositions.
    \end{enumerate}
    Which of the above statements are \textbf{correct}?
    \begin{enumerate}[label=(\alph*)]
        \item (i) and (iii) only
        \item (ii) and (iv) only
        \item (i), (ii), and (iii) only
        \item (i), (iii), and (iv) only
    \end{enumerate}

    \item Which of the following statements about the Svedberg's unit ('S') in ribosomes is \textbf{correct}?
    \begin{enumerate}[label=(\alph*)]
        \item It is an additive unit of mass, which is why $50\text{S} + 30\text{S}$ equals $80\text{S}$
        \item It stands for the sedimentation coefficient and is indirectly a measure of density and size
        \item It measures the number of amino acids present in the ribosomal proteins
        \item It is applicable only to membrane-bound organelles
    \end{enumerate}

    \item \textbf{Statement I:} In eukaryotic cells, $80\text{S}$ ribosomes can be found attached to the outer surface of the endoplasmic reticulum as well as freely in the cytoplasm.\\
    \textbf{Statement II:} The endoplasmic reticulum bearing ribosomes on its surface is called rough endoplasmic reticulum (RER) and is frequently observed in cells actively involved in protein synthesis and secretion.\\
    Choose the correct option:
    \begin{enumerate}[label=(\alph*)]
        \item Both Statement I and Statement II are correct
        \item Both Statement I and Statement II are incorrect
        \item Statement I is correct but Statement II is incorrect
        \item Statement I is incorrect but Statement II is correct
    \end{enumerate}

    \item Identify the \textbf{incorrect} comparison between a prokaryotic cell and a chloroplast:
    \begin{enumerate}[label=(\alph*)]
        \item Both possess $70\text{S}$ ribosomes made of $50\text{S}$ and $30\text{S}$ subunits
        \item Both contain circular DNA molecules
        \item Both possess a cell wall made of peptidoglycan
        \item Both carry out protein synthesis independently using their own ribosomes
    \end{enumerate}

    \item Which of the following statements is \textbf{true} regarding the DNA found inside mitochondria and chloroplasts?
    \begin{enumerate}[label=(\alph*)]
        \item Mitochondria contain a single circular DNA molecule, whereas chloroplasts contain small, double-stranded circular DNA molecules
        \item Both mitochondria and chloroplasts contain linear DNA associated with histone proteins
        \item Neither mitochondria nor chloroplasts contain RNA molecules
        \item Mitochondrial DNA is single-stranded, while chloroplast DNA is triple-stranded
    \end{enumerate}

    \item \textbf{Assertion (A):} Chloroplasts and mitochondria are often referred to as semi-autonomous organelles.\\
    \textbf{Reason (R):} Both organelles possess their own DNA, RNA, and $70\text{S}$ ribosomes, enabling them to synthesize some of their own proteins.\\
    Choose the correct option:
    \begin{enumerate}[label=(\alph*)]
        \item Both (A) and (R) are true and (R) is the correct explanation of (A)
        \item Both (A) and (R) are true but (R) is not the correct explanation of (A)
        \item (A) is true but (R) is false
        \item Both (A) and (R) are false
    \end{enumerate}

    \item Match Column I with Column II regarding structures and their exact locations:
    \begin{center}
    \begin{tabular}{@{}ll@{}}
        \toprule
        \textbf{Column I (Structure)} & \textbf{Column II (Location)} \\
        \midrule
        A. Cristae & 1. Stroma of chloroplast \\
        B. Stroma lamellae & 2. Inner membrane of mitochondrion \\
        C. $70\text{S}$ ribosomes & 3. Cytoplasm of eukaryotic cells \\
        D. $80\text{S}$ ribosomes & 4. Mitochondrial matrix and chloroplast stroma \\
        \bottomrule
    \end{tabular}
    \end{center}
    Choose the correct option:
    \begin{enumerate}[label=(\alph*)]
        \item A-2, B-1, C-4, D-3
        \item A-1, B-2, C-4, D-3
        \item A-2, B-4, C-1, D-3
        \item A-3, B-1, C-4, D-2
    \end{enumerate}

    \item A student examines a cell under an electron microscope and observes an organelle bounded by two membranes, containing flattened membranous sacs arranged in stacks like piles of coins, and flat tubules connecting these stacks. Which pigment(s) would be expected inside this organelle?
    \begin{enumerate}[label=(\alph*)]
        \item Only xanthophylls and carotene, with complete absence of chlorophyll
        \item Chlorophyll and carotenoid pigments
        \item Only water-soluble anthocyanins
        \item No pigments, as it is a leucoplast
    \end{enumerate}

    \item \textbf{Statement I:} All plastids are interchangeable and share a common evolutionary origin, possessing similar internal structure.\\
    \textbf{Statement II:} Leucoplasts are colourless plastids that store reserve food materials in the form of starch, oils/fats, or proteins.\\
    Choose the correct option:
    \begin{enumerate}[label=(\alph*)]
        \item Both Statement I and Statement II are correct
        \item Both Statement I and Statement II are incorrect
        \item Statement I is correct but Statement II is incorrect
        \item Statement I is incorrect but Statement II is correct
    \end{enumerate}

    \item Which of the following cell types is expected to have the \textbf{highest} number of mitochondria per cell?
    \begin{enumerate}[label=(\alph*)]
        \item A dormant seed embryo cell
        \item A metabolically inactive epidermal cell
        \item An actively contracting flight muscle cell of a bird
        \item A mature sieve tube element
    \end{enumerate}

    \item Identify the \textbf{incorrect} statement about the stroma lamellae of chloroplasts:
    \begin{enumerate}[label=(\alph*)]
        \item They are flat membranous tubules
        \item They connect the thylakoids of the different grana
        \item They are located inside the mitochondrial matrix
        \item They are part of the internal membrane system within the chloroplast stroma
    \end{enumerate}

    \item \textbf{Assertion (A):} Ribosomes are the only organelles found in both prokaryotic and eukaryotic cells.\\
    \textbf{Reason (R):} Prokaryotic cells lack all membrane-bound organelles, and ribosomes are non-membrane bound granular structures.\\
    Choose the correct option:
    \begin{enumerate}[label=(\alph*)]
        \item Both (A) and (R) are true and (R) is the correct explanation of (A)
        \item Both (A) and (R) are true but (R) is not the correct explanation of (A)
        \item (A) is true but (R) is false
        \item Both (A) and (R) are false
    \end{enumerate}

    \item Consider the following pairs of organelles and their staining/visibility characteristics:
    \begin{enumerate}[label=(\roman*)]
        \item Mitochondria — Easily visible under light microscope without staining
        \item Plastids — Easily observed under the microscope as they are large
        \item Ribosomes — First observed under the electron microscope as dense particles
        \item Nucleolus — Non-membrane bound structure present in the nucleoplasm
    \end{enumerate}
    Which of the above pairs is \textbf{incorrectly} matched?
    \begin{enumerate}[label=(\alph*)]
        \item (i) only
        \item (ii) only
        \item (iii) and (iv)
        \item (i) and (ii)
    \end{enumerate}

    \item Which of the following sequences accurately represents the arrangement from the outermost boundary to the innermost core of a chloroplast?
    \begin{enumerate}[label=(\alph*)]
        \item Outer membrane $\rightarrow$ Intermembrane space $\rightarrow$ Inner membrane $\rightarrow$ Stroma $\rightarrow$ Thylakoid membrane $\rightarrow$ Thylakoid lumen
        \item Outer membrane $\rightarrow$ Inner membrane $\rightarrow$ Thylakoid lumen $\rightarrow$ Thylakoid membrane $\rightarrow$ Stroma
        \item Inner membrane $\rightarrow$ Outer membrane $\rightarrow$ Stroma $\rightarrow$ Grana $\rightarrow$ Cristae
        \item Outer membrane $\rightarrow$ Stroma $\rightarrow$ Inner membrane $\rightarrow$ Thylakoid lumen $\rightarrow$ Thylakoid membrane
    \end{enumerate}

    \item \textbf{Statement I:} In a eukaryotic cell, both the outer nuclear membrane and the rough endoplasmic reticulum bear $80\text{S}$ ribosomes on their cytoplasmic surfaces.\\
    \textbf{Statement II:} The smooth endoplasmic reticulum (SER) lacks ribosomes and is the major site for synthesis of lipid.\\
    Choose the correct option:
    \begin{enumerate}[label=(\alph*)]
        \item Both Statement I and Statement II are correct
        \item Both Statement I and Statement II are incorrect
        \item Statement I is correct but Statement II is incorrect
        \item Statement I is incorrect but Statement II is correct
    \end{enumerate}

    \item How many of the following structures possess their own $70\text{S}$ ribosomes?
    \begin{center}
        \textit{Mitochondria, Chloroplast, Bacterium, Blue-green algae, Rough Endoplasmic Reticulum, Mycoplasma}
    \end{center}
    \begin{enumerate}[label=(\alph*)]
        \item Three
        \item Four
        \item Five
        \item Six
    \end{enumerate}

    \item Which of the following statements about plastids in Euglenoids is \textbf{true} according to NCERT?
    \begin{enumerate}[label=(\alph*)]
        \item Euglenoids completely lack plastids and are strictly heterotrophic
        \item Plastids are found in all plant cells and in euglenoids
        \item Euglenoids possess only leucoplasts and lack photosynthetic pigments
        \item Euglenoid plastids have peptidoglycan walls
    \end{enumerate}

    \item Match Column I with Column II regarding key discoveries and terms in cell biology:
    \begin{center}
    \begin{tabular}{@{}ll@{}}
        \toprule
        \textbf{Column I (Term / Scientist)} & \textbf{Column II (Associated Feature)} \\
        \midrule
        A. George Palade (1953) & 1. Chain of ribosomes attached to a single mRNA \\
        B. Polysome & 2. Observed ribosomes as dense particles under electron microscope \\
        C. Cristae & 3. Stacks of thylakoids resembling piles of coins \\
        D. Grana & 4. Infoldings of inner mitochondrial membrane \\
        \bottomrule
    \end{tabular}
    \end{center}
    Choose the correct option:
    \begin{enumerate}[label=(\alph*)]
        \item A-2, B-1, C-4, D-3
        \item A-1, B-2, C-4, D-3
        \item A-2, B-4, C-1, D-3
        \item A-4, B-1, C-2, D-3
    \end{enumerate}

\end{enumerate}

\newpage

\section*{\color{primary} Part II: Quick Answer Key}

\begin{center}
\begin{tabular}{@{}ccccc@{}}
    \toprule
    \textbf{Q. No. / Ans} & \textbf{Q. No. / Ans} & \textbf{Q. No. / Ans} & \textbf{Q. No. / Ans} & \textbf{Q. No. / Ans} \\
    \midrule
    \textbf{1.} (c) & \textbf{10.} (a) & \textbf{19.} (b) & \textbf{28.} (a) & \textbf{37.} (c) \\
    \textbf{2.} (c) & \textbf{11.} (b) & \textbf{20.} (c) & \textbf{29.} (b) & \textbf{38.} (c) \\
    \textbf{3.} (c) & \textbf{12.} (d) & \textbf{21.} (a) & \textbf{30.} (a) & \textbf{39.} (a) \\
    \textbf{4.} (c) & \textbf{13.} (d) & \textbf{22.} (b) & \textbf{31.} (c) & \textbf{40.} (a) \\
    \textbf{5.} (a) & \textbf{14.} (c) & \textbf{23.} (a) & \textbf{32.} (a) & \textbf{41.} (a) \\
    \textbf{6.} (b) & \textbf{15.} (c) & \textbf{24.} (c) & \textbf{33.} (a) & \textbf{42.} (a) \\
    \textbf{7.} (c) & \textbf{16.} (b) & \textbf{25.} (b) & \textbf{34.} (a) & \textbf{43.} (c) \\
    \textbf{8.} (d) & \textbf{17.} (b) & \textbf{26.} (a) & \textbf{35.} (b) & \textbf{44.} (b) \\
    \textbf{9.} (b) & \textbf{18.} (a) & \textbf{27.} (a) & \textbf{36.} (a) & \textbf{45.} (a) \\
    \bottomrule
\end{tabular}
\end{center}

\vspace{12pt}

\section*{\color{primary} Part III: Detailed Solutions \& NCERT Explanations}

\begin{enumerate}[label=\textbf{\arabic*.}, leftmargin=*, itemsep=6pt]

    \item \textbf{(c)} NCERT Class XI (Chapter 8) explicitly states that mitochondria, unless specifically stained, are \textbf{not} easily visible under the microscope. Their dimensions are $0.2\text{--}1.0\,\mu\text{m}$ (average $0.5\,\mu\text{m}$) in diameter and $1.0\text{--}4.1\,\mu\text{m}$ in length.

    \item \textbf{(c)} Statements (ii), (iii), and (iv) are correct. Statement (i) is incorrect because the number of mitochondria per cell is \textbf{variable} depending on the physiological activity of the cells.

    \item \textbf{(c)} Assertion (A) is true: the inner membrane forms infoldings called cristae towards the matrix. Reason (R) is false because the cristae \textbf{increase} (not decrease) the surface area.

    \item \textbf{(c)} The mitochondrial matrix possesses a single circular DNA molecule, a few RNA molecules, and \textbf{$70\text{S}$ ribosomes} (not $80\text{S}$ ribosomes), along with components required for protein synthesis.

    \item \textbf{(a)} Outer membrane forms the continuous limiting boundary; inner membrane divides the lumen into two aqueous compartments; cristae are infoldings towards the matrix increasing surface area; mitochondria are sites of aerobic respiration producing ATP.

    \item \textbf{(b)} According to NCERT, mitochondria divide by \textbf{fission} (binary fission).

    \item \textbf{(c)} Statement I is correct. Statement II is incorrect because mitochondria are the sites of \textbf{aerobic respiration} (not anaerobic fermentation), producing cellular energy in the form of ATP, hence called 'power houses' of the cell.

    \item \textbf{(d)} All four statements are verbatim facts from NCERT Class XI (Plastids section). Plastids are found in all plant cells and in euglenoids, are large enough to be easily observed under the microscope, and are classified into chloroplasts, chromoplasts, and leucoplasts.

    \item \textbf{(b)} NCERT states that in chromoplasts, \textbf{fat-soluble carotenoid pigments like carotene, xanthophylls and others} are present, giving the part of the plant a yellow, orange, or red colour.

    \item \textbf{(a)} Leucoplasts are colourless plastids divided into: \textbf{Amyloplasts} (store carbohydrates/starch, e.g., potato), \textbf{Elaioplasts} (store oils and fats), and \textbf{Aleuroplasts} (store proteins).

    \item \textbf{(b)} Both statements are true according to NCERT, as chloroplasts are double membrane-bound and the inner chloroplast membrane is relatively less permeable than the outer membrane, but (R) is simply a general structural fact and not the mechanistic explanation of (A).

    \item \textbf{(d)} NCERT states the exact opposite: \textit{Chlamydomonas} (a green alga) has \textbf{1 chloroplast per cell}, whereas a mesophyll cell has \textbf{$20\text{--}40$ chloroplasts per cell}.

    \item \textbf{(d)} All four statements accurately describe the internal organization of a chloroplast as given in NCERT.

    \item \textbf{(c)} Statement I is correct (thylakoid membrane encloses a lumen). Statement II is incorrect because chlorophyll pigments are present in the \textbf{thylakoids}, not in the stroma.

    \item \textbf{(c)} The stroma of the chloroplast contains enzymes required for the synthesis of carbohydrates and proteins, small double-stranded circular DNA molecules, and $70\text{S}$ ribosomes.

    \item \textbf{(b)} NCERT defines the endomembrane system (ER, Golgi complex, lysosomes, vacuoles) as membranous organelles whose functions are coordinated. Since the functions of mitochondria, chloroplast, and peroxisomes are \textbf{not coordinated} with these components, they are not considered part of the endomembrane system.

    \item \textbf{(b)} Eukaryotic cytoplasmic ribosomes are \textbf{$80\text{S}$}, whereas the ribosomes found inside chloroplasts and mitochondria are smaller and are of the \textbf{$70\text{S}$} type.

    \item \textbf{(a)} Ribosomes are non-membrane bound organelles present in both prokaryotes and eukaryotes. In eukaryotes, they are found in the cytoplasm ($80\text{S}$), on the rough ER ($80\text{S}$), and inside mitochondria and chloroplasts ($70\text{S}$).

    \item \textbf{(b)} Ribosomes are the granular structures first observed under the electron microscope as dense particles by \textbf{George Palade (1953)}.

    \item \textbf{(c)} Statement I is correct. Statement II is incorrect because 'S' (Svedberg's unit) stands for the sedimentation coefficient; it is \textbf{indirectly} a measure of density and size (not a direct measure of molecular weight).

    \item \textbf{(a)} Eukaryotic ($80\text{S}$) ribosomes consist of two subunits: \textbf{$60\text{S}$ and $40\text{S}$}. Prokaryotic ($70\text{S}$) ribosomes consist of \textbf{$50\text{S}$ and $30\text{S}$} subunits.

    \item \textbf{(b)} NCERT Chapter 8 (Prokaryotic Cells) specifies that prokaryotic ribosomes are associated with the plasma membrane and are about \textbf{$15\text{ nm}$ by $20\text{ nm}$} in size.

    \item \textbf{(a)} Several ribosomes may attach to a single mRNA and form a chain called polyribosomes or polysome, which translate the mRNA into proteins.

    \item \textbf{(c)} Both mitochondria and chloroplasts are double membrane-bound organelles containing their own circular DNA and $70\text{S}$ ribosomes. Thylakoids/grana and carotenoids are unique to chloroplasts, while cristae are unique to mitochondria.

    \item \textbf{(b)} Chromoplasts contain \textbf{fat-soluble} carotenoid pigments (carotene, xanthophylls), not water-soluble pigments.

    \item \textbf{(a)} Mitochondrion diameter: $0.2\text{--}1.0\,\mu\text{m}$ (avg $0.5\,\mu\text{m}$); Chloroplast length: $5\text{--}10\,\mu\text{m}$; Prokaryotic ribosome: $15\text{ nm} \times 20\text{ nm}$; Chloroplast width: $2\text{--}4\,\mu\text{m}$.

    \item \textbf{(a)} Carbohydrate synthesis (dark reaction/Calvin cycle) occurs in the chloroplast stroma precisely because the stroma houses all the necessary enzymes for carbohydrate and protein synthesis.

    \item \textbf{(a)} Statements (i) and (iii) are correct. Statement (ii) is wrong because the \textbf{inner} compartment (not outer) is filled with the matrix. Statement (iv) is wrong because both membranes have their own \textbf{specific} enzymes.

    \item \textbf{(b)} In ribosomes, 'S' (Svedberg's unit) stands for the sedimentation coefficient, which is indirectly a measure of density and size (hence Svedberg values are non-additive).

    \item \textbf{(a)} Both statements accurately reflect NCERT descriptions of ribosomes and the rough endoplasmic reticulum (RER).

    \item \textbf{(c)} Chloroplasts do not have a peptidoglycan cell wall; they are bounded by a double lipid-bilayer membrane. However, they share $70\text{S}$ ribosomes and circular DNA with prokaryotes.

    \item \textbf{(a)} NCERT specifically states that the mitochondrial matrix possesses a \textbf{single circular DNA molecule}, whereas the chloroplast stroma contains \textbf{small, double-stranded circular DNA molecules}.

    \item \textbf{(a)} Because mitochondria and chloroplasts possess their own circular DNA, RNA, and $70\text{S}$ ribosomes, they can replicate by fission and synthesize some of their own proteins, making them semi-autonomous.

    \item \textbf{(a)} Cristae are on the inner mitochondrial membrane; stroma lamellae are in the chloroplast stroma; $70\text{S}$ ribosomes are in the mitochondrial matrix and chloroplast stroma; $80\text{S}$ ribosomes are in the eukaryotic cytoplasm.

    \item \textbf{(b)} The description matches a \textbf{chloroplast} (stacks of coins = grana/thylakoids; flat tubules = stroma lamellae). Chloroplasts contain chlorophyll and carotenoid pigments.

    \item \textbf{(a)} Plastids share a common origin and can interconvert (e.g., chloroplasts turning into chromoplasts during fruit ripening), and leucoplasts are colourless nutrient-storing plastids.

    \item \textbf{(c)} The number of mitochondria per cell correlates directly with physiological/metabolic activity. An actively contracting flight muscle cell requires immense ATP and thus has a very high number of mitochondria.

    \item \textbf{(c)} Stroma lamellae are flat membranous tubules connecting thylakoids of different grana inside the \textbf{chloroplast stroma}, not the mitochondrial matrix.

    \item \textbf{(a)} Prokaryotes lack all membrane-bound organelles, leaving non-membrane bound ribosomes as the only organelle common to both prokaryotic and eukaryotic cells.

    \item \textbf{(a)} Pair (i) is incorrect because mitochondria are \textbf{not} easily visible under the light microscope unless specifically stained.

    \item \textbf{(a)} Moving inward in a chloroplast: Outer membrane $\rightarrow$ Intermembrane space $\rightarrow$ Inner membrane $\rightarrow$ Stroma $\rightarrow$ Thylakoid membrane $\rightarrow$ Thylakoid lumen.

    \item \textbf{(a)} Both statements are accurate according to NCERT Class XI (Endomembrane System).

    \item \textbf{(c)} Five structures possess $70\text{S}$ ribosomes: Mitochondria, Chloroplast, Bacterium, Blue-green algae (cyanobacteria), and Mycoplasma (all prokaryotes have $70\text{S}$ ribosomes). Rough Endoplasmic Reticulum bears eukaryotic $80\text{S}$ ribosomes.

    \item \textbf{(b)} NCERT explicitly states in the opening sentence on plastids: \textit{"Plastids are found in all plant cells and in euglenoids."}

    \item \textbf{(a)} George Palade (1953) discovered ribosomes; Polysome is a chain of ribosomes on a single mRNA; Cristae are inner mitochondrial membrane infoldings; Grana are stacks of thylakoids resembling piles of coins.

\end{enumerate}

\end{document}